% Part: proof-theory
% Chapter: normalization
% Section: normalization-thm

\documentclass[../../../include/open-logic-section]{subfiles}

\begin{document}

\olfileid{pt}{nor}{thm}

\olsection{د عادي کولو قضیه}

!!^a{proof} په \emph{عادي} بڼه کښې دے که د پرېکون
برخې نۀ لري۔ په څرګند ډول دا هغه وخت، او يوازې هغه وخت،
داسې وي چې نۀ پرې د ځاے بدلولو بدلون او نۀ د راکمولو
بدلون تطبيقولے شو۔ د !!a{proof} \emph{عادي کولو} معنا
دا ده چې د ځاے بدلولو او راکمولو بدلونونه پرې تطبيق
کړو، تر څو هېڅ پرېکون پاتې نۀ شي، او په عادي ثبوت ئې
واړوو۔ دا چې دا کار تل شونے دے، د \emph{عادي کولو قضیه} ده۔

\begin{thm}\ollabel{thm:normalization}
هر \Log{N1i}-!!{proof} $\delta$ عادي کېږي، يعنې داسې
!!a{proof}~$\delta'$ ته اوړېدے شي چې پرېکونونه نۀ لري۔
\end{thm}

\begin{proof}
  د~$\delta$ د پرېکون پر رتبه او د پرېکون پر اوږدوالي
  استقرا کوو۔ د استقرا فرضيه دا ده چې هر !!{proof} چې
  يا د پرېکون ټيټه رتبه لري، يا په هماغه رتبه کښې د
  پرېکون کم اوږدوالے لري، عادي کېږي۔
  \textbf{د سرچينه‌يي اندازې وضاحت OLCMP-170:} دلته
  دويمه اندازه د پرېکون اوږدوالے، يعنې د تر ټولو لوړې
  رتبې د پرېکون برخو د اوږدوالو جمع ده؛ د ټول ثبوت
  اوږدوالے نۀ دے۔ د سرچينې لنډ ``اوږدوالے'' د همدغه
  مخکيني تعريف او د استقرا د پاتې بيان سره څرګند شو۔

  که د پرېکون اوږدوالے~$0$ وي، پرېکونونه نۀ شته او
  $\delta$ له مخکې عادي دے۔ نو فرض کړئ چې د~$\delta$ د
  پرېکون رتبه او اوږدوالے~$> 0$ دي۔ د تر ټولو پاسنۍ او
  ښي خوا ته د تر ټولو لوړې رتبې د پرېکون برخه غوره
  کړئ۔ که د هغې اوږدوالے $>1$ وي، د ځاے بدلولو بدلون
  تطبيق کړئ۔ که اوږدوالے ئې~$1$ وي، اړوند د راکمولو
  بدلون تطبيق کړئ۔ په دې يو نوے !!{proof} جوړېږي چې يا
  د پرېکون هماغه رتبه او کم اوږدوالے لري، يا د پرېکون
  ټيټه رتبه لري؛ نو د استقرا فرضيه پرې تطبيقېږي۔
\end{proof}

د پورته ثبوت تګلاره، چې تل د تر ټولو پاسني او ښي
خوا ته د تر ټولو لوړې رتبې پرېکون سره کار کوي، داسې
ځانګړے ترتيب راښيي چې بدلونونه په هغې کښې پر
!!a{proof} تطبيق شي او عادي !!{proof} ورکړي۔ خو د
حيرانتيا خبره دا ده چې په حقيقت کښې ترتيب اهميت نۀ
لري۔ د ځاے بدلولو او راکمولو بدلونونه په هر ترتيب
تطبيقول په پاے کښې يو عادي ثبوت ورکوي۔

د برخې او پرېکون تعريف د \Log{N2i} دپاره هم په
آسانۍ جوړېږي، او د \Log{N1i} د ځاے بدلولو او راکمولو
بدلونونه نېغ \Log{N2i} ته ژباړل کېږي۔ نو د عادي کولو
قضیه د \Log{N2i} دپاره هم صحيح ده۔

\begin{thm}\ollabel{thm:normalization-N2i}
هر \Log{N2i}-!!{proof} $\delta$ عادي کېږي، يعنې داسې
!!a{proof}~$\delta'$ ته اوړېدے شي چې پرېکونونه نۀ لري۔
\end{thm}

\end{document}
